\section{Background and Preliminaries}
\label{sec:prelim}

In this section, we establish the formal mathematical definitions, notation, and objective functions that govern modern temporal foundation models and large language model adaptations.

\subsection{Formal Problem Definitions}

\subsubsection{Temporal Sequences}
A univariate time series of length $T$ is defined as an ordered sequence of real-valued scalar observations:
\begin{equation}
    \mathbf{x} = [x_1, x_2, \dots, x_T]^\top \in \mathbb{R}^T,
\end{equation}
where $x_t \in \mathbb{R}$ represents the measurement recorded at discrete timestamp $t \in \{1, \dots, T\}$.

A multivariate time series comprises $C$ co-evolving channels or variates observed synchronously across $T$ timestamps:
\begin{equation}
    \mathbf{X} = [\mathbf{x}^{(1)}, \mathbf{x}^{(2)}, \dots, \mathbf{x}^{(C)}]^\top \in \mathbb{R}^{C \times T},
\end{equation}
where $\mathbf{X}_{:, t} \in \mathbb{R}^C$ is the vector of all variate measurements at timestamp $t$.

\subsubsection{Channel Modeling Paradigms}
When processing multivariate series $\mathbf{X}$, architectures adopt one of two primary structural paradigms:
\begin{itemize}[leftmargin=*]
    \item \emph{Channel-Independence (CI)}: Introduced prominently in PatchTST~\cite{nie2023patchtst} and Timer~\cite{liu2024timer}, CI decomposes the multivariate series into $C$ distinct univariate sequences:
    \begin{equation}
        \mathbf{X} \mapsto \{\mathbf{x}^{(1)}, \mathbf{x}^{(2)}, \dots, \mathbf{x}^{(C)}\},
    \end{equation}
    which are fed into the shared model backbone independently with tied parameters. CI eliminates the risk of overfitting spurious cross-channel correlations and natively handles variable channel counts $C$ at inference.
    \item \emph{Channel-Mixing (CM)}: Models such as MOIRAI~\cite{woo2024moirai} and TTM~\cite{ekambaram2024ttm} employ full cross-variate attention or specialized channel-mixing MLPs:
    \begin{equation}
        \mathbf{H} = \text{Attn}_{\text{channel}}(\mathbf{X}) \in \mathbb{R}^{C \times T \times D},
    \end{equation}
    explicitly capturing physical inter-sensor dependencies at the cost of higher computational complexity.
\end{itemize}

\subsection{Tokenization and Input Representation}

Because continuous temporal values lack an inherent lexicon, models deploy specialized mapping operators $\mathcal{T}: \mathbb{R}^* \to \mathbb{R}^D$ to construct transformer-compatible token embeddings.

\subsubsection{Subseries Patching}
Rather than embedding individual scalar time-steps (which leads to quadratic attention overhead $\mathcal{O}(T^2)$ and severe information sparsity), the majority of modern TSFMs (e.g., TimesFM~\cite{das2024timesfm}, Timer~\cite{liu2024timer}, MOMENT~\cite{goswami2024moment}) partition a sequence $\mathbf{x}$ into $N$ contiguous subseries patches of length $P$ with stride $S$:
\begin{equation}
    N = \left\lfloor \frac{T - P}{S} \right\rfloor + 1.
\end{equation}
The $i$-th patch $\mathbf{p}_i = [x_{(i-1)S + 1}, \dots, x_{(i-1)S + P}]^\top \in \mathbb{R}^P$ is mapped to a $D$-dimensional latent vector $\mathbf{e}_i \in \mathbb{R}^D$ via an affine projection:
\begin{equation}
    \mathbf{e}_i = \mathbf{p}_i \mathbf{W}_p + \mathbf{b}_p + \mathbf{e}_{\text{pos}, i},
\end{equation}
where $\mathbf{W}_p \in \mathbb{R}^{P \times D}$, $\mathbf{b}_p \in \mathbb{R}^D$, and $\mathbf{e}_{\text{pos}, i} \in \mathbb{R}^D$ denotes a learnable or sinusoidal positional encoding. Patching preserves local temporal shape dynamics and compresses sequence length from $T$ to $N \ll T$.

\subsubsection{Scalar Quantization and Categorical Binning}
In contrast to continuous linear patching, Amazon Chronos~\cite{ansari2024chronos} introduces a pure language-analog tokenization. Given an input sequence $\mathbf{x}$, it computes instance-level scale statistics $(\mu, \sigma)$ via mean scaling:
\begin{equation}
    \widetilde{x}_t = \frac{x_t - \mu}{\sigma + \epsilon}, \quad \mu = \frac{1}{T}\sum_{t=1}^T x_t, \quad \sigma = \frac{1}{T}\sum_{t=1}^T |x_t - \mu|.
\end{equation}
The normalized scalars $\widetilde{x}_t$ are discretized into $K$ uniform bins over $[-B, B]$:
\begin{equation}
    q_t = \text{clip}\left(\left\lfloor \frac{\widetilde{x}_t - (-B)}{2B / K} \right\rfloor, 0, K - 1\right),
\end{equation}
where $q_t \in \{0, 1, \dots, K-1\}$ acts as a discrete categorical token id, enabling standard cross-entropy language modeling.

\subsection{Downstream Task Formulations}

\subsubsection{Zero-Shot and Supervised Forecasting}
Given an historical context window $\mathbf{X}_{1:T} \in \mathbb{R}^{C \times T}$, the forecasting objective is to predict the future sequence of length $H$:
\begin{equation}
    \widehat{\mathbf{Y}} = \widehat{\mathbf{X}}_{T+1:T+H} \in \mathbb{R}^{C \times H}.
\end{equation}
In \emph{probabilistic forecasting}, the model estimates the conditional predictive distribution:
\begin{equation}
    p(\mathbf{Y} \mid \mathbf{X}_{1:T}) = \prod_{h=1}^H p(y_{T+h} \mid \mathbf{X}_{1:T}, y_{T+1:T+h-1}).
\end{equation}
This distribution is parameterized either parametrically (e.g., Student's $t$ parameters in Lag-Llama~\cite{rasul2023lagllama} or Gaussian mixtures in MOIRAI~\cite{woo2024moirai}) or non-parametrically through categorical bin logits (Chronos~\cite{ansari2024chronos}).

\subsection{Pre-training Objectives}

\subsubsection{Autoregressive Generation (Next-Patch / Next-Token)}
Decoder-only foundation models (TimesFM~\cite{das2024timesfm}, Timer~\cite{liu2024timer}, Chronos~\cite{ansari2024chronos}) optimize an autoregressive likelihood:
\begin{equation}
    \mathcal{L}_{\text{AR}}(\Theta) = -\sum_{i=1}^N \log p_\Theta(\mathbf{p}_{i+1} \mid \mathbf{p}_1, \dots, \mathbf{p}_i).
\end{equation}
For continuous regression heads (Timer, TimesFM), this is computed using Mean Squared Error (MSE) or Mean Absolute Error (MAE):
\begin{equation}
    \mathcal{L}_{\text{reg}}(\Theta) = \frac{1}{N} \sum_{i=1}^N \|\mathbf{p}_{i+1} - \widehat{\mathbf{p}}_{i+1}\|_2^2.
\end{equation}

\subsubsection{Masked Autoencoding (MAE)}
Encoder-based foundation models (MOMENT~\cite{goswami2024moment}, PatchTST~\cite{nie2023patchtst}) apply a binary masking vector $\mathbf{m} \in \{0, 1\}^N$ over patches, optimizing patch reconstruction over the masked set $\mathcal{M} = \{i \mid m_i = 1\}$:
\begin{equation}
    \mathcal{L}_{\text{Mask}}(\Theta) = \frac{1}{|\mathcal{M}|} \sum_{i \in \mathcal{M}} \|\mathbf{p}_i - f_\Theta(\mathbf{X}_{\backslash \mathcal{M}})\|_2^2.
\end{equation}
Masked autoencoding yields bidirectional temporal representations particularly suited for downstream classification, anomaly detection, and imputation.
